\documentclass[10pt,a4paper]{article}
\usepackage[margin=1.9cm]{geometry}
\usepackage{amsmath,amssymb,graphicx,booktabs,array,longtable,hyperref}
\setlength{\parskip}{4pt}\setlength{\parindent}{0pt}
\newcommand{\xb}{\bar x}
\newcommand{\dens}{{\rm dens}}
\newcommand{\ion}{{\rm ion}}
\title{An analytic toy for the shape of $D_\ell^{\rm cross}(x_{\rm HII})$\\\large kSZ$^2\times$galaxy cross-power during reionisation: what sets the curve?}
\author{S.\ Upadhye \;--\; \texttt{ksz-lae-xcorr}, \texttt{correlation/toy\_model.py}, \texttt{scripts/50\_toy\_model.py}}
\date{29 October 2026 (v2: symbols defined, notation cleaned)}
\begin{document}\maketitle

\section{Question}
La Plante, Sipple \& Lidz (2022) Fig.~7 plots $D_\ell^{\rm cross}$, the cross-power of a filtered, squared kSZ map with a galaxy map from a redshift window of width $\Delta z=1$, against the mean ionised fraction $x_{\rm HII}$ of that window. It rises from $x_{\rm HII}\simeq0$, peaks near $x_{\rm HII}\simeq 0.1$--$0.2$ and has fallen to almost nothing by $x_{\rm HII}\simeq0.8$. Is that shape already explained by a few lines of algebra, and which ingredient makes it?

\section{Symbols (every symbol, once)}
Notation rules used throughout: the \emph{pixel} index is always $j$; the words ``dens'' and ``ion'' (never an index) label the density and ionisation effects; $Q$ is the slice power; $\chi$ is a distance; $G$ is the growth factor; $D_\ell$ is the angular power; $R$ is not used.

{\small
\begin{longtable}{>{$}p{2.4cm}<{$} p{14.0cm}}\toprule
\multicolumn{2}{l}{\textbf{Maps and the estimator}}\\\midrule
\Theta(\hat n) & the kSZ map: the kSZ temperature fluctuation (units of $T_{\rm CMB}$) summed along the line of sight from all pixels at $z\ge6$, $\Theta=\sum_j\theta_j$.\\
\theta_j & kSZ contribution of line-of-sight pixel $j$ (1 Mpc long): $\theta_j=-W_j\,(1+\delta_j)\,x_{e,j}\,v_j/c$.\\
W_j & kSZ weight of pixel $j$: $W_j=K(z_j)\,\Delta\chi_j$, the optical depth per pixel, damped by $e^{-\tau}$ (see $K$).\\
\delta_j,\ x_{e,j},\ v_j & matter overdensity, free-electron fraction and line-of-sight peculiar velocity at pixel $j$; $c$ = speed of light.\\
F(\ell) & the CMB filter, a function of multipole. It keeps the multipoles where the kSZ is not drowned by the primary CMB and noise. $F\Theta$ means: the map $\Theta$ with this filter applied in Fourier space.\\
S & the \emph{squared filtered map}, $S=[F\Theta]^2$. After filtering only small scales remain, so $S$ is the local small-scale kSZ power: bright where the small-scale kSZ is strong.\\
g(\hat n) & the galaxy (tracer) map of one redshift window: $g=\int dz\,w(z)\,b_g(z)\,\delta_m$, a mean-zero overdensity map.\\
w(z) & window weight, a top hat of width $\Delta z=1$ with $\int w\,dz=1$.\\
C_\ell^{Sg},\ D_\ell & $C_\ell^{Sg}$ is the angular cross-power of $S$ with $g$; $D_\ell\equiv D_\ell^{\rm cross}=\ell(\ell+1)C_\ell^{Sg}/2\pi$ (times $T_{\rm CMB}^2$ for $\mu$K$^2$; in the toy this is absorbed in $A$).\\\midrule
\multicolumn{2}{l}{\textbf{Functions of redshift (all ordinary numbers at each $z$)}}\\\midrule
z,\ \chi(z) & redshift; comoving distance to $z$. Multipole and wavenumber are linked by $k=\ell/\chi$.\\
\bar x(z) & mean ionised fraction of the universe at $z$ (the reionisation history). $x_{\rm HII}$ of a window is $\bar x$ averaged over it.\\
\tau(z),\ K(z) & optical depth to $z$, $\tau=\int\tau_{\rm pref}\,\bar x\,(1+z)^2\,d\chi$; and $K=\tau_{\rm pref}(1+z)^2e^{-\tau}$, the kSZ weight per unit length ($\tau_{\rm pref}=n_{H,0}\sigma_T$).\\
G(z) & linear growth factor of matter fluctuations, normalised to $G(0)=1$.\\
v(z) & growth scaling of the peculiar velocity, $v=E\,\Omega_m(z)^{0.55}G/(1+z)$ ($=aHfG$ in units of $H_0$; $E=H/H_0$). Only its redshift dependence is used.\\
b_g(z) & tracer bias: $\delta_{\rm tracer}\simeq b_g\,\delta_m$ on large scales. $b_g=2.1(1+z)-5.3$ (BlueTides fit) for every tracer.\\
P_{\rm lin}(k) & linear matter power spectrum at $z=0$ (Eisenstein--Hu), so $G^2P_{\rm lin}$ is the power at $z$.\\
p(z) & how fast reionisation is going at $z$: $p=-d\bar x/dz\ge0$.\\
b_{zm}(k) & zreion bias, $b_0/(1+k/k_0)^{0.2}$, $b_0=1/1.686$, $k_0=0.9\,h$Mpc$^{-1}$. $z_{\rm mean}=8$ is the reference redshift of the zreion model and $z_{\rm re}(\mathbf x)$ is the redshift at which the point $\mathbf x$ reionises: $z_{\rm re}=z_{\rm mean}+(1+z_{\rm mean})\,b_{zm}\,\delta_L(z_{\rm mean})$.\\\midrule
\multicolumn{2}{l}{\textbf{The toy's own quantities}}\\\midrule
Q(z) & small-scale kSZ$^2$ power of the slice of sky at $z$, \emph{per unit redshift}, at the filter scale (this replaces the symbol $s$ of v1, which clashed with a distance).\\
\delta_L & the long-wavelength (large-scale) matter overdensity of that slice: the part the galaxies trace.\\
Q_{\delta}(z) & the response of $Q$ to $\delta_L$: $Q_\delta=\partial Q/\partial\delta_L$, i.e.\ $Q\to Q+Q_\delta\,\delta_L$. It has the same units as $Q$ and is an \textbf{ordinary function of $z$, not an operator}. $Q_\delta=Q_\dens+Q_\ion$.\\
Q_\dens,\ Q_\ion & the two parts of $Q_\delta$: the effect of density (more matter, more electrons) and the effect of ionisation (denser regions reionise earlier). Both are ordinary numbers at each $z$.\\
\epsilon & ratio of patchy-bubble power to smooth-ionisation power at the filter scale. Fitted ($\simeq8$).\\
r_v & small optional correction from the velocity variance; set to $0$.\\
A & overall amplitude per $\ell$ (fitted; absorbs $T_{\rm CMB}^2$, the filter's share of the power and the velocity normalisation).\\
\Delta z,\ z_0 & width ($=1$) and centre of the galaxy window.\\\bottomrule
\end{longtable}}

\section{The model, in four steps}
\textbf{Step 1: power of one slice.} With $K$, $v$ and $\bar x$ as above,
\begin{equation}
 Q(z)=K^2v^2\left[\xb^2+\epsilon\,\xb(1-\xb)\right].
\end{equation}
The first term is the kSZ of a smoothly ionised slice with ionised fraction $\xb$. The second is the extra power of the bubbles: a binary ionisation field has variance $\xb(1-\xb)$, which is zero when the slice is fully neutral or fully ionised and largest when half ionised. $\epsilon$ sets its size and is the only shape parameter.

\textbf{Step 2: how $Q$ responds to the large-scale density $\delta_L$.}
\begin{align}
 Q_\dens &= (2+r_v)\,Q, \\
 Q_\ion &= \frac{\partial Q}{\partial \xb}\;\frac{\partial \xb}{\partial\delta_L}, \qquad \frac{\partial Q}{\partial \xb}=K^2v^2\left[2\xb+\epsilon(1-2\xb)\right], \\
 \frac{\partial \xb}{\partial\delta_L} &=(1+z_{\rm mean})\,b_{zm}(k)\,\frac{G(z_{\rm mean})}{G(z)}\;p(z).
\end{align}
In words: $Q_\dens$ is the first line because the power goes as $(1+\delta)^2$, which grows by $2\delta_L$ when the slice is denser ($r_v=0$). $Q_\ion$ is the change of $Q$ with $\xb$ times the change of $\xb$ with $\delta_L$; the latter is zreion: a denser region reionises earlier ($z_{\rm re}$ shifts up), so at redshift $z$ its ionised fraction is higher by the shift times $p$.
$Q_\dens$ and $Q_\ion$ are plain numbers at each $z$, added together; nothing is applied to anything. The sign of $Q_\ion$ is that of $2\xb+\epsilon(1-2\xb)$, which is positive for $\xb<\epsilon/[2(\epsilon-1)]\simeq0.57$ and \emph{negative} above it.

\textbf{Step 3: cross-power with the galaxies.} The galaxies trace $\delta_L$ with bias $b_g$. In the Limber approximation only the slice at the same redshift as the galaxies correlates with them (Kramer et al.\ 2025, Eqs.~9--10, with a tracer in place of $\tau$):
\begin{equation}
 D_\ell=\frac{\ell(\ell+1)}{2\pi}\,\frac{A}{\Delta z}\int_{z_0-\Delta z/2}^{z_0+\Delta z/2}\!dz\;\big(Q_\dens+Q_\ion\big)\,b_g(z)\,\frac{G(z)^2\,P_{\rm lin}(\ell/\chi)}{\chi^2}.
\end{equation}
The power is taken per unit redshift; the factor $d\chi/dz$ is omitted (including it leaves the best $\epsilon$ unchanged, misfit $0.33\to0.43$).

\textbf{Step 4: the reionisation history.} La Plante's digitised fiducial $\bar x(z)$, made monotonic and Gaussian-smoothed ($\sigma_z=0.3$) so that $p$ is usable, with $\bar x=1$ below the table and an $e^{-(z-11.86)}$ tail above its end ($\bar x=0.06$).

\begin{table}[h]\centering\small
\begin{tabular}{r|rrrrrr}\toprule
$z$ & 6 & 8 & 9 & 10 & 11 & 12\\\midrule
$\xb$ & 0.91 & 0.44 & 0.24 & 0.13 & 0.08 & 0.05\\
$\tau$ & 0.046 & 0.062 & 0.066 & 0.068 & 0.070 & 0.071\\
$K^2v^2\;[10^{-10}]$ & 0.42 & 0.87 & 1.2 & 1.6 & 2.0 & 2.6\\
$p=-d\xb/dz$ & 0.10 & 0.26 & 0.15 & 0.07 & 0.03 & 0.04\\
$b_g$ & 9.4 & 13.6 & 15.7 & 17.8 & 19.9 & 22.0\\
$G(z)$ & 0.18 & 0.14 & 0.13 & 0.12 & 0.11 & 0.10\\\bottomrule
\end{tabular}
\\[2pt]{\small Redshift functions entering the integrand (fiducial history).}
\end{table}

\section{Fit}
For each $\epsilon$ and each $\ell\in\{500,1000,3000\}$ the amplitude $A$ is the log-space least-squares value against La Plante's digitised Fig.~7 curve over $x_{\rm HII}=0.1$--$0.9$ (40 points); the misfit is the rms of $\ln D_{\rm LP}-\ln(A\,D_{\rm toy})$. $\epsilon$ is the value minimising the misfit averaged over the three $\ell$.

\begin{center}\includegraphics[width=\linewidth]{toy_vs_laplante.pdf}\end{center}
Orange: La Plante (fiducial). Black: toy ($\epsilon=8.25$), with its density part ($Q_\dens$, dashed) and ionisation part ($Q_\ion$, dotted). Right: misfit versus $\epsilon$. The plot axis stops at zero, so the dotted curve appears to end near $x_{\rm HII}\simeq0.58$; in fact it continues \emph{below zero} there.

\section{Results}
\begin{itemize}\setlength\itemsep{1pt}
\item \textbf{Best $\epsilon\simeq8$} (misfit 0.28/0.34/0.36 in log at $\ell=500/1000/3000$; $A=7.0,3.4,1.7\times10^{6}$). Peak at $x_{\rm HII}\simeq0.24$ (La Plante: 0.12--0.19); falls to $\sim0$ by $x_{\rm HII}\simeq0.8$. The shape is therefore reproduced to within the digitisation and the toy's crudeness; the peak sits slightly too late.
\item For $\epsilon\lesssim3$ (smooth term competitive) the misfit is $\gtrsim0.7$: the peak moves to $x_{\rm HII}\simeq0.3$ ($\epsilon=3$), 0.55 ($\epsilon=1$), 0.7 ($\epsilon=0$) and the curve keeps 30--70\% of its peak value at $x_{\rm HII}=0.9$, whereas La Plante keeps $\sim2\%$ (toy at $\epsilon=8$: 6\%).
\item The two parts at $\ell=1000$, relative to the peak of the total ($\epsilon=8.25$):
\begin{center}\small\begin{tabular}{r|rrrrrr}\toprule
$x_{\rm HII}$ & 0.05 & 0.13 & 0.24 & 0.45 & 0.72 & 0.91\\\midrule
$Q_\dens$ part & 0.19 & 0.32 & 0.43 & 0.49 & 0.37 & 0.20\\
$Q_\ion$ part & 0.44 & 0.49 & 0.57 & 0.26 & $-0.20$ & $-0.14$\\
total & 0.63 & 0.81 & 1.00 & 0.75 & 0.17 & 0.06\\\bottomrule
\end{tabular}\end{center}
The density part alone peaks near $x_{\rm HII}\simeq0.45$ and cannot fall fast enough. The ionisation part peaks near 0.2 and turns \emph{negative} above $\xb\simeq0.57$: a dense region is then already almost fully ionised, so it has \emph{fewer} bubbles than an average region and less small-scale power. This cancels most of the positive density part, and together with $p\to0$ as reionisation ends produces the late fall.
\item Robustness: $z_{\rm mean}=7,8,9$ give the same $\epsilon$ and misfit, because $(1+z_{\rm mean})G(z_{\rm mean})$ is nearly constant. The minimum in $\epsilon$ is sharp (misfit 0.33 at 8.25, 0.50 at 10, 1.0 at 11, $\ge1.3$ for $\epsilon\ge12$) because the total depends on this cancellation; $\epsilon$ should not be quoted as a measurement.
\end{itemize}

\section{What it says}
The early rise and the late fall of $D_\ell^{\rm cross}(x_{\rm HII})$ are produced by \emph{patchy} ionisation: the kSZ$^2$ power of a partly ionised slice (growing as $\xb(1-\xb)$) is correlated with the large-scale density mainly through where bubbles form (the zreion response, $\propto p(z)$), not through the smooth density term. This is consistent with the component study, where $v_{\rm los}$ and the ionisation morphology both follow the kSZ curve. A tracer that is a better proxy for $\delta_L$ at the patchy epoch (larger effective $b_g$, less noise) therefore carries more of this signal.

\section{Caveats and checks to do}
\begin{itemize}\setlength\itemsep{1pt}
\item Squeezed-limit, slice-diagonal toy: no cross-slice pairs, no line-of-sight cancellation of the smooth term. $\epsilon\sim8$ may partly stand in for that cancellation, so it is a statement about the relative size of the patchy and smooth terms in the \emph{estimator}, not a measured physical ratio.
\item $x_{\rm HII}<0.1$ is controlled by the first points of the digitised history (a small wiggle remains); $b_g(z)$ is one shape for all tracers; only the shape is fitted, the amplitude is not predicted.
\item Not yet done: the same overlay on our own measured curves (\texttt{--ours-root}) and a direct test of $\epsilon$ from the simulation (power of the $x_e$-only versus $v$-only kSZ$^2$ maps per window, scripts 48/49).
\end{itemize}

\section{Reproduce}
\begin{quote}\small\begin{verbatim}
python scripts/50_toy_model.py --out-dir data/plots/2026-10-23_toy \
   [--ours-root data/plots/2026-10-16_tracer_cross] [--eps 8.25]
\end{verbatim}\end{quote}
\end{document}
